\documentclass[10pt,a4paper]{amsart}
\usepackage[margin=19mm]{geometry}
\usepackage{amsmath,amssymb,mathtools}
\usepackage[scr=rsfs,cal=euler]{mathalfa}
\usepackage{xfrac}
\usepackage[unicode,hidelinks,bookmarks=false]{hyperref}
\newcommand{\ie}{i.e.\ }
\newcommand{\cf}{cf.\ }
\newcommand{\resp}{respectively\ }
\newcommand{\Subsection}{\subsection}
\providecommand{\nobreakdash}{\nobreak\hbox{-}}
\setlength{\emergencystretch}{2em}
\multlinegap0pt
\usepackage{bm}
\usepackage{bbm}
\usepackage{dsfont}
\usepackage{xspace}
\usepackage{cancel}
\usepackage{mathtools}
\usepackage{amsthm}
\makeatletter
\@ifundefined{theorem}{\newtheorem{theorem}{Theorem}}{}
\@ifundefined{definition}{\newtheorem{definition}[theorem]{Definition}}{}
\@ifundefined{lemma}{\newtheorem{lemma}[theorem]{Lemma}}{}
\makeatother
\title[A characterization of ball quotient stacks]{A characterization of ball quotient stacks}
\author{Chirantan Chowdhury, Matteo Costantini, Aryaman Patel}
\date{Source: arXiv:2608.04688v1 (CC BY 4.0)}
\begin{document}
\maketitle

\noindent\emph{Source and licence.} A characterization of ball quotient stacks. Version arXiv:2608.04688v1 (CC BY 4.0), distributed under the Creative Commons Attribution 4.0 International licence, as recorded in the arXiv licence metadata for this exact version and shown on the abstract page https://arxiv.org/abs/2608.04688v1. Full rights notice: licenses/arxiv-2608.04688-CC-BY-4.0.txt.\par\smallskip

\noindent\textit{Editorial scope.} Counted unit: the Lemma whose source label is \texttt{Completeness-lem} in the article (source0.tex lines 454--458, source label \texttt{Completeness-lem}), delivered together with the article's own complete original proof of it (source0.tex lines 459--479, one \texttt{proof} environment of 21 lines). No mathematics is written, repaired or re-derived: statement and proof are the article's own text line for line. Required context retained uncounted: Uniformizng VHS (lines 202--206). Every definition, display and label of those lines is retained unchanged. Omitted from this package and not redistributed: 1--201, 207--453, 480--711. Nothing inside a retained excerpt is omitted or altered: every retained body is delivered exactly as the article's lines are written, so no omission receipt of any kind accompanies this package. Citations: the retained excerpts carry no citation command at all, so no bibliography entry of the article is transcribed and no reference is redistributed. Reference adaptation: none was needed and none was made; every reference inside the delivered unit resolves inside it, so no unresolved reference or citation is produced. Printed slips kept literal: no printed word, symbol, hypothesis or number of the counted statement or of its proof was altered. Layout and class: the article's own class is \texttt{amsart} with packages including fontenc, amsmath, amssymb, amsthm, mathrsfs, mathtools, graphicx, varioref, appendix, enumitem, xspace, ifthen, comment, xy; the wrapper is the lane's pinned \texttt{amsart} editorial wrapper at 10pt on A4, which restates the theorem-like environments the retained text needs and adds the article's own macro definitions that the retained text needs (none), with extra wrapper packages (bm, bbm, dsfont, xspace, cancel, mathtools). The delivered numbering is the unit's own. Assets: no retained span contains a figure environment, a graphics-inclusion command, a PDF-page inclusion, a table or a TikZ picture; no image file is copied into this package, the package declares no asset dependency and no image file is redistributed with it. Licence: version arXiv:2608.04688v1 of this article is distributed under the Creative Commons Attribution 4.0 International licence, as recorded in the arXiv licence metadata for this exact version and shown on the abstract page https://arxiv.org/abs/2608.04688v1. A full rights notice is delivered at licenses/arxiv-2608.04688-CC-BY-4.0.txt. Characters: every retained excerpt is pure ASCII, so the delivered engine, XeLaTeX with its default Latin Modern text font, reports no missing character.
\begin{definition}\label{Uniformizng VHS}
A \emph{uniformizing system of log Hodge bundles} on a log pair $(\mathscr{X},\mathscr{D})$ is a principal system of log Hodge bundles $(P,\tau)$ such that the morphism $\tau:\mathcal{T}_{\mathscr{X}}(-\log \mathscr{D})\to P\times_K\mathfrak{g}^{-1,1}$ is an isomorphism. 

Moreover, if $(P|_{\mathscr{X}-\mathscr{D}},\tau|_{\mathscr{X}-\mathscr{D}})$ admits a metric reduction $P_H$ with structure group $K_0$ such that the curvature of the connection $D_H$ on the $G_0$-bundle $P_H\times_{K_0}G_0$ is zero, then the triple $(P|_{\mathscr{X}-\mathscr{D}},\tau|_{\mathscr{X}-\mathscr{D}},P_H)$ is a \emph{uniformizing variation of Hodge structure} on $\mathscr{X}-\mathscr{D}$.
\end{definition}

\begin{lemma}\label{Completeness-lem}
Let $(\mathscr{X},\mathscr{D})$ be a proper log pair and let $(P,\tau)$ be a uniformizing system of log Hodge bundles on $(\mathscr{X},\mathscr{D})$ for a Hodge group $G_0$ of Hermitian type, as in Definition \ref{Uniformizng VHS}. Suppose that there is a metric reduction $P_H\subset P|_{\mathscr{X}-\mathscr{D}}$ such that $(P|_{\mathscr{X}-\mathscr{D}},\tau|_{\mathscr{X}-\mathscr{D}},P_H)$ is a uniformizing variation of Hodge structure on $\mathscr{X}-\mathscr{D}$.

If the metric induced by $P_H$ on $\mathcal{T}_{\mathscr{X}-\mathscr{D}}$ is complete, then there is an isomorphism $\widetilde{\mathscr{X}-\mathscr{D}}\xrightarrow{\sim}G_0/K_0=:\mathcal{D}$.
\end{lemma}

\begin{proof}
Let $(\widetilde{P},\widetilde{\tau},\widetilde{P}_H)$ denote the pullback of the uniformizing variation of Hodge structure $(P|_{\mathscr{X}-\mathscr{D}},\tau|_{\mathscr{X}-\mathscr{D}},P_H)$ to the universal covering stack $\widetilde{\mathscr{X}-\mathscr{D}}$.

The flat connection $D_H$ on the bundle $P_H\times_{K_0}G_0$ (see Definition \ref{Uniformizng VHS}) pulls back to a flat connection on the bundle $\widetilde{P}_H\times_{K_0}G_0$ and induces a trivialization $\widetilde{P}_H\times_{K_0}G_0\cong(\widetilde{\mathscr{X}-\mathscr{D}})\times G_0$. Consider the composition
\[
\widetilde{P}_H\hookrightarrow\widetilde{P}_H\times_{K_0}G_0\xrightarrow{\sim}(\widetilde{\mathscr{X}-\mathscr{D}})\times G_0\to G_0
\]
which is a $K_0$-equivariant morphism. This induces a morphism 
\[
\phi:\widetilde{\mathscr{X}-\mathscr{D}}\to G_0/K_0=:\mathcal{D}.
\] 
By construction, we have $\widetilde{P}_H\cong\phi^*G_0$, where $G_0$ is viewed as a principal $K_0$-bundle on $\mathcal{D}$.

The metric $h_H$ on $P|_{\mathscr{X}-\mathscr{D}}\times_K\mathfrak{g}^{-1,1}$ pulls back to a metric $\widetilde{h}_H$ on $\widetilde{P}\times_K\mathfrak{g}^{-1,1}$. Note that this is the same as the metric induced by the reduction $\widetilde{P}_H\subset\widetilde{P}$. Since $\widetilde{\tau}:\mathcal{T}_{\widetilde{\mathscr{X}-\mathscr{D}}}\to\widetilde{P}\times_K\mathfrak{g}^{-1,1}$ is a isomorphism by Definition \ref{Uniformizng VHS}, this may be viewed as a metric on $\mathcal{T}_{\widetilde{\mathscr{X}-\mathscr{D}}}$.

On $\mathcal{D}$, we have an isomorphism $\mathcal{T}_{\mathcal{D}}\cong G_0\times_{K_0}\mathfrak{g}^{-1,1}$, and since $\widetilde{P}_H\cong\phi^*G_0$, it follows that $\mathcal{T}_{\widetilde{\mathscr{X}-\mathscr{D}}}\cong\phi^*\mathcal{T}_{\mathcal{D}}$. This implies that $\phi$ is \'etale and since $\mathcal{D}$ is an analytic space, so is $\widetilde{\mathscr{X}-\mathscr{D}}$.

Moreover, we have $\widetilde{h}_H=\phi^*h_\mathcal{D}$, where $h_{\mathcal{D}}$ is the canonical Hermitian metric on $\mathcal{T}_{\mathcal{D}}$. Thus $\phi$ is in fact a local isometry and $\widetilde{h}_H$ in invariant under the $\pi_1(\mathscr{X}-\mathscr{D})$-action on $\widetilde{\mathscr{X}-\mathscr{D}}$.

Therefore, if $h_H$ is a complete metric, then so is $\widetilde{h}_H$ and $\phi:\widetilde{\mathscr{X}-\mathscr{D}}\to\mathcal{D}$ is a covering map. Since $\widetilde{\mathscr{X}-\mathscr{D}}$ and $\mathcal{D}$ are both simply connected analytic spaces, it follows that $\phi$ is an isomorphism, as desired. 
\end{proof}

\end{document}
